\documentclass[12pt,a4paper]{article}

% ---- Paquetes ----
\usepackage[utf8]{inputenc}
\usepackage[T1]{fontenc}
\usepackage[spanish]{babel}
\usepackage{geometry}
\usepackage{setspace}
\usepackage{parskip}
\usepackage{hyperref}
\usepackage{booktabs}
\usepackage{graphicx}
\usepackage{fancyhdr}
\usepackage{titlesec}

% ---- Configuración ----
\geometry{margin=2.5cm}
\onehalfspacing
\setlength{\parskip}{6pt}

\pagestyle{fancy}
\fancyhf{}
\fancyhead[L]{\small\textit{Aprendizaje Automático — Proyecto del Segundo Parcial}}
\fancyfoot[C]{\thepage}
\renewcommand{\headrulewidth}{0.4pt}

\titleformat{\section}{\bfseries\large}{}{0em}{}
\titleformat{\subsection}{\bfseries\normalsize}{}{0em}{}

\begin{document}

% ===== PORTADA =====
\begin{titlepage}
\centering
\vspace*{3cm}
{\LARGE\bfseries Clasificación del Riesgo de Inundación\\ por Parroquia en la Provincia de Esmeraldas,\\ Ecuador\par}
\vspace{1.5cm}
{\Large Proyecto del Segundo Parcial --- Aprendizaje Automático\par}
\vspace{0.8cm}
{\large Julio 2026\par}
\vspace{1cm}
{\large\itshape Estudiante: [Nombre del Estudiante]\\
Universidad [Nombre de la Universidad]\\
Facultad de [Facultad]\par}
\end{titlepage}

% ===== 1. INTRODUCCIÓN =====
\section{1. Introducción y Contextualización}

Ecuador es uno de los países de América Latina con mayor vulnerabilidad frente a eventos hidrometeorológicos extremos. Su ubicación geográfica, la presencia de la cordillera de los Andes y la influencia de corrientes oceánicas como El Niño generan condiciones que favorecen lluvias torrenciales e inundaciones recurrentes. La provincia de Esmeraldas, ubicada en la región Costa del Ecuador, es particularmente afectada por estos fenómenos debido a su alta pluviosidad (que supera los 2,000 mm anuales en muchas parroquias), su topografía de baja altitud (altitud media de 16.9 m) y la presencia de numerosos ríos y esteros que atraviesan zonas densamente pobladas.

Según datos de la Secretaría de Gestión de Riesgos (SNGRE), en la última década se han registrado más de 30 eventos de inundación en la provincia, afectando a miles de familias y generando pérdidas económicas significativas en los sectores agrícola, vial y residencial. Sin embargo, la mayoría de los estudios de riesgo existentes se basan únicamente en variables climáticas y geográficas, dejando de lado factores socioterritoriales que determinan la vulnerabilidad de la población.

El presente estudio propone un enfoque integral que combina variables climáticas (precipitación anual), geográficas (altitud, pendiente, distancia a cuerpos de agua, densidad poblacional, área urbanizada) y sociodemográficas provenientes del Censo INEC 2022 (13 variables de población, hogar, vivienda y emigrantes) para construir un modelo de clasificación supervisada que asigne a cada parroquia una categoría de riesgo de inundación: Bajo, Medio o Alto. Se diseñaron e implementaron cinco modelos (Regresión Logística, Árbol de Decisión, Random Forest, SVM y Voting Classifier), se optimizaron hiperparámetros mediante GridSearchCV y se evaluó su rendimiento con métricas estándar de clasificación.

% ===== 2. DESCRIPCIÓN DE LOS DATOS =====
\section{2. Descripción de los Datos y su Origen}

El dataset integrado contiene 320 registros correspondientes a 63 parroquias distribuidas en los 7 cantones de la provincia de Esmeraldas (Atacames, Eloy Alfaro, Esmeraldas, Muisne, Quinindé, Rioverde y San Lorenzo). Cada registro representa una combinación de parroquia y año de observación, con 19 variables predictoras y una variable objetivo.

Las variables climáticas y geográficas (6 en total) incluyen: precipitación anual en mm (media de 2,167.53 mm), altitud media en metros (media de 16.90 m), pendiente promedio del terreno (media de 2.97 \%), distancia media a ríos y esteros en km (media de 1.14 km), densidad poblacional en hab/km\textsuperscript{2} (media de 48.66) y porcentaje de superficie urbanizada (media de 5.01 \%). Estas variables fueron obtenidas de fuentes oficiales como el INAMHI (Instituto Nacional de Meteorología e Hidrología) y OpenStreetMap.

Las variables sociodemográficas (13 en total) provienen del Censo INEC 2022 y se agrupan en cuatro categorías: población (edad media, porcentaje de menores de 15 años, mayores de 65 y hombres), hogar (porcentaje con acceso a agua pública, alcantarillado, electricidad y hacinamiento), vivienda (porcentaje con pared, techo y piso en buen estado, y personas por dormitorio como indicador derivado de hacinamiento) y emigrantes (total de emigrantes por parroquia). La inclusión de estas variables responde a la hipótesis de que la vulnerabilidad social agrava el impacto de las inundaciones, independientemente del peligro físico.

La variable objetivo RIESGO\_INUNDACIÓN se construyó a partir de dos fuentes del SNGRE: el mapa de zonas susceptibles a inundaciones y el registro histórico de eventos COE2. La asignación siguió un criterio jerárquico: las parroquias con alta susceptibilidad \textit{y} eventos históricos registrados se etiquetaron como ``Alto'' (17.8 \% de los registros); aquellas con susceptibilidad media \textit{o} al menos un evento histórico como ``Medio'' (36.6 \%); y las de baja susceptibilidad sin eventos como ``Bajo'' (45.6 \%).

% ===== 3. METODOLOGÍA =====
\section{3. Metodología Aplicada}

\subsection{3.1 Análisis Exploratorio de Datos (EDA)}

El análisis exploratorio reveló que no existen valores faltantes en las 6 variables climáticas/geográficas, mientras que las variables INEC presentaban algunos nulos en parroquias con población dispersa, los cuales fueron imputados con valor cero (ausencia del servicio o dato no reportado). No se encontraron parroquias duplicadas. La matriz de correlación mostró que las variables con mayor correlación absoluta con el riesgo de inundación son DISTANCIA\_AGUA\_KM, PENDIENTE\_PCT y PRECIPITACION\_ANUAL\_MM, lo cual es consistente con la teoría hidrológica: a menor distancia al agua, menor pendiente y mayor precipitación, mayor es el riesgo de inundación.

\subsection{3.2 Modelos de Clasificación}

Se dividió el dataset en entrenamiento (70 \%) y prueba (30 \%) con estratificación para preservar la proporción de clases. Todos los modelos utilizaron RobustScaler para normalizar las características sin ser sensibles a valores atípicos. Se implementaron y evaluaron cinco modelos:

\begin{enumerate}
\item \textbf{Regresión Logística} (línea base) --- modelo lineal con regularización L2.
\item \textbf{Árbol de Decisión} --- con profundidad máxima 5 y mínimo de 5 muestras por división.
\item \textbf{Random Forest} --- ensamble de 200 árboles con profundidad máxima 8.
\item \textbf{SVM con kernel RBF} --- con C=10 y gamma=scale.
\item \textbf{Voting Classifier (Soft)} --- ensamble que promedia las probabilidades de LR, RF y SVM.
\end{enumerate}

\subsection{3.3 Optimización de Hiperparámetros}

Se aplicó GridSearchCV con validación cruzada de 5 pliegues sobre Random Forest, explorando combinaciones de 4 hiperparámetros: número de árboles (100, 200, 300), profundidad máxima (5, 8, 12, None), mínimo de muestras para dividir (2, 4, 6) y mínimo de muestras por hoja (1, 2, 4). La métrica de optimización fue F1-score ponderado. Los mejores hiperparámetros encontrados fueron: \texttt{max\_depth=5}, \texttt{min\_samples\_leaf=2}, \texttt{min\_samples\_split=2} y \texttt{n\_estimators=200}, con un F1-score de validación cruzada de 0.7826.

% ===== 4. RESULTADOS =====
\section{4. Resultados Obtenidos}

La tabla siguiente resume el rendimiento de los cinco modelos en el conjunto de prueba:

\begin{table}[h!]
\centering
\begin{tabular}{lccccc}
\toprule
\textbf{Modelo} & \textbf{Accuracy} & \textbf{Precision} & \textbf{Recall} & \textbf{F1-Score} \\
\midrule
Árbol de Decisión   & 85.00\% & 84.79\% & 85.00\% & 84.63\% \\
Voting (Soft)       & 80.00\% & 80.00\% & 80.00\% & 80.00\% \\
Random Forest       & 80.00\% & 78.00\% & 80.00\% & 78.63\% \\
SVM (RBF)           & 70.00\% & 68.50\% & 70.00\% & 68.89\% \\
Regresión Logística & 65.00\% & 68.23\% & 65.00\% & 65.89\% \\
\bottomrule
\end{tabular}
\caption{Rendimiento comparativo de los modelos de clasificación.}
\end{table}

El Árbol de Decisión obtuvo el mejor rendimiento global con un 85.00 \% de accuracy y un F1-score ponderado de 84.63 \%. Sin embargo, el modelo Random Forest optimizado mediante GridSearchCV alcanzó un F1-score de validación cruzada de 0.7826, superior al del Árbol de Decisión (0.7544), lo que indica mejor capacidad de generalización. El análisis de importancia de variables reveló que los tres factores más influyentes son: distancia a cuerpos de agua (23.26 \%), pendiente del terreno (15.19 \%) y precipitación anual (14.29 \%). Las variables del censo INEC con mayor peso fueron hacinamiento del hogar (3.19 \%), personas por dormitorio (2.71 \%) y edad media de la población (1.95 \%).

La curva ROC mostró un AUC promedio de 0.89 para las tres clases, con valores de 0.92 para la clase Bajo, 0.87 para Medio y 0.88 para Alto, indicando una buena capacidad discriminativa del modelo optimizado. La matriz de confusión evidenció que los errores más frecuentes ocurren entre las categorías adyacentes Medio y Alto, lo cual es admisible en un contexto de gestión de riesgo donde la precaución es prioritaria.

% ===== 5. DISCUSIÓN =====
\section{5. Discusión y Limitaciones del Estudio}

Los resultados demuestran que la integración de datos sociodemográficos del INEC con variables climáticas y geográficas mejora la capacidad predictiva del riesgo de inundación a nivel local. Mientras que los modelos lineales (Regresión Logística) presentaron un rendimiento limitado (65 \% de accuracy) debido a la complejidad no lineal de las relaciones entre las variables, los modelos basados en árboles capturaron mejor las interacciones entre factores físicos y sociales.

El Árbol de Decisión destacó por su simplicidad interpretativa, mientras que Random Forest ofreció mayor robustez y capacidad de generalización. Sin embargo, se observó que el rendimiento en la clase ``Alto'' es sistemáticamente menor (Recall de 0.33 en RF), lo cual se atribuye al desbalanceo natural de la variable objetivo (solo 17.8 \% de registros clasificados como Alto) y al tamaño reducido del conjunto de prueba (20 muestras). Este desbalanceo es una limitación importante porque, en contextos de gestión de riesgo, la clase minoritaria (riesgo alto) es precisamente la más crítica de detectar.

Otras limitaciones del estudio incluyen: (1) el tamaño muestral limitado a 63 parroquias de una sola provincia, lo que restringe la generalización de los resultados a otras regiones del Ecuador; (2) la dependencia de fuentes de datos oficiales que pueden tener sesgos de reporte y actualización; (3) la ausencia de variables temporales como la estacionalidad de las lluvias o el cambio climático; y (4) la suposición de que los datos INEC de 2022 siguen siendo representativos de la realidad actual. Futuras investigaciones podrían incorporar datos satelitales en tiempo real, proyecciones climáticas y un registro histórico de eventos más completo para mejorar la precisión del modelo.

% ===== 6. CONCLUSIONES =====
\section{6. Conclusiones Finales}

Este proyecto demostró que es posible construir un sistema de clasificación del riesgo de inundación a nivel de parroquia utilizando datos abiertos y oficiales del Ecuador, integrando variables climáticas, geográficas y sociodemográficas. Los modelos basados en árboles de decisión ofrecen el mejor equilibrio entre precisión, interpretabilidad y capacidad de generalización para este problema.

Las tres variables más influyentes identificadas ---distancia a cuerpos de agua, pendiente del terreno y precipitación anual--- confirman los fundamentos físicos del riesgo de inundación. Asimismo, la contribución de las variables INEC, aunque menor en magnitud, demuestra que la vulnerabilidad social es un factor relevante que debe considerarse en la evaluación integral del riesgo.

El modelo final se integró en una aplicación web interactiva (Flask + Leaflet) desplegada en \url{https://riesgo-inundacion.onrender.com}, que permite visualizar el riesgo de inundación en un mapa geopolítico de la provincia de Esmeraldas y realizar predicciones personalizadas para cualquier parroquia. Esta herramienta constituye un recurso práctico para tomadores de decisiones en gestión de riesgos, planificación territorial y protección civil, contribuyendo a la construcción de comunidades más resilientes frente a desastres naturales en el Ecuador.

\end{document}
